\documentclass{article}
\usepackage{amsmath,amssymb}
\usepackage{xurl}
\usepackage[hidelinks]{hyperref}
% Shared commands for notes in docs/paper-gaps/.

\DeclareUnicodeCharacter{2080}{\ifmmode{}_{0}\else\textsubscript{0}\fi}
\DeclareUnicodeCharacter{2081}{\ifmmode{}_{1}\else\textsubscript{1}\fi}
\DeclareUnicodeCharacter{2082}{\ifmmode{}_{2}\else\textsubscript{2}\fi}
\DeclareUnicodeCharacter{2099}{\ifmmode{}_{n}\else\textsubscript{n}\fi}

\newcommand{\Lean}{\textsc{Lean}}
\ExplSyntaxOn
\NewDocumentCommand{\leanid}{m}{
  \ifmmode
    \textsf{\detokenize{#1}}
  \else
    \group_begin:
    \urlstyle{sf}
    \tl_set:Nn \l_tmpa_tl {#1}
    \tl_replace_all:Nnn \l_tmpa_tl {\_} {_}
    \exp_args:NV \path \l_tmpa_tl
    \group_end:
  \fi
}
\ExplSyntaxOff
\providecommand{\tr}{\operatorname{tr}}
\newcommand{\tnleanIssueBase}{https://github.com/LionSR/TNLean/issues/}
\newcommand{\tnleanPrBase}{https://github.com/LionSR/TNLean/pull/}
\newcommand{\tnleanDocBase}{https://sirui-lu.com/TNLean/docs/}
\newcommand{\ghissue}[1]{\href{\tnleanIssueBase#1}{GitHub issue \##1}}
\newcommand{\ghpr}[1]{\href{\tnleanPrBase#1}{PR \##1}}
\newcommand{\doclink}[1]{\href{\tnleanDocBase#1.html}{\path{#1}}}

% Dirac notation and the identity operator, as in blueprint/src/macros/common.tex.
% \providecommand keeps note-local definitions authoritative.
\usepackage{dsfont}
\providecommand{\ket}[1]{|#1\rangle}
\providecommand{\bra}[1]{\langle #1|}
\providecommand{\braket}[2]{\langle #1 | #2 \rangle}
\providecommand{\ketbra}[2]{|#1\rangle\!\langle #2|}
\providecommand{\Id}{\mathds{1}}
\providecommand{\one}{\mathds{1}}
\providecommand{\C}{\mathbb{C}}
\providecommand{\MN}[1]{M_{#1}(\C)}
\providecommand{\MD}{M_D(\C)}

% Verdict marker: \gapnote{<kind>}{<status>}. Kind is the policy
% classification (clarification, local-correction, scope-restriction,
% unfaithful, false-source, open-gap); status is open, wip (elimination
% underway), resolved, or historical. Machine-read by the site index;
% prints a verdict line.
\newcommand{\gapnote}[2]{\par\noindent\textbf{Verdict:} #1 (#2).\par}

\title{Repeated Overlapping Blocks with a Small Weight Sum}
\author{}
\date{2026-10-02}
\begin{document}
\maketitle
\gapnote{scope-restriction}{resolved}

\section{What the source states}
Let $A^i=\bigoplus_{j=1}^b\operatorname{diag}(\mu_{j,1},\dots,\mu_{j,m_j})
    \otimes A_j^i$ with normal $A_j$ \cite[eq.~(S2)]{Malz2024Preparation}. The source
assumes ``Without loss of generality we assume $| \mu_{j,k} | \le 1$ with at least one
of them having magnitude exactly one'' \cite[after eq.~(S2)]{Malz2024Preparation}, puts
$\beta_j=\sum_k\mu_{j,k}^N$ \cite[eq.~(S4)]{Malz2024Preparation}, and bounds the error
$\varepsilon=1-\lvert\braket{\widetilde\phi_N}{\phi_N}\rvert$ of its approximating state
by $O\bigl((N/q)e^{-\gamma q/\xi_{\mathrm{diag}}}\bigr)$
\cite[Lemma~1$'$(ii)]{Malz2024Preparation}, with no factor depending on the weights.

\section{The former restriction}
For blocks with multiplicities whose $q$-site states may overlap, the bound for the corrected
approximating state of \cite{gap:mswc24_repeated_block_corrected_state}, at the rate of
\cite{gap:mswc24_block_form_mixed_overlap}, was first proved for complex weights with
$\lvert\mu_{j,k}\rvert\leq1$ in the form
\begin{align*}
    \varepsilon\leq\frac{Cy\,e^{Cy}}{\min(1,b)^{1/2}},\qquad y=Me^{-\gamma q/\xi},\qquad
    b=\sum_j\lvert\beta_j\rvert^2,
\end{align*}
and without the factor only for real weights $0\leq\mu_{j,k}\leq1$ of which one equals one,
where $b\geq1$. Within the source's normalization $b<1$ can occur: for copy weights $(1,-1)$
at odd $N$ the sum $\beta_j$ of that block vanishes, and $b$ is then the contribution of the
other blocks, which may be exponentially small in $N$. The factor came from the proof: the
overlaps $z_j=\braket{\Omega_j}{\phi_M(L_j^\dagger P)}$ of the rows of the positive part $P$
with the fixed-point tensors of the blocks were estimated in absolute terms,
$\lvert z_j-\beta_j\rvert\leq C_jy\,e^{C_jy}$, and dividing by $b^{1/2}$ to normalize lost the
factor.

\section{The relative estimate}
The factor is removed by estimating the positive part relative to the weights of the blocks.
With the copy norms $c_j=(\sum_k\lvert\mu_{j,k}\rvert^{2q})^{1/2}$ and the copy isometries
$L_j$ of \cite{gap:mswc24_repeated_block_corrected_state}, the blocked map is
$B=\widetilde BC$ with $\widetilde B=\sum_jB_jL_j^\dagger$ and $C=\sum_jc_jL_jL_j^\dagger$. Put
$C^{-1}=\sum_jc_j^{-1}L_jL_j^\dagger$, $P_\infty=\sum_jc_jL_j(\sqrt{\sigma_j}^{\,T}\otimes
\mathbf 1)L_j^\dagger$ and $\Delta=\widetilde B^\dagger\widetilde B-\sum_jL_j(\sigma_j^T\otimes
\mathbf 1)L_j^\dagger$. Then $X=(P-P_\infty)C^{-1}$ solves the Sylvester equation
\begin{align*}
    PX+XP_\infty=C\Delta\Pi,\qquad\Pi=\sum_jL_jL_j^\dagger .
\end{align*}
If $\lambda>0$ bounds every $\sigma_j^T\otimes\mathbf 1$ from below and
$\lVert\Delta\rVert\leq\lambda/2$, then $P^2=C\widetilde B^\dagger\widetilde BC\geq(\lambda/2)C^2$.
In eigenbases $P=\sum_ap_af_af_a^\dagger$ and $P_\infty=\sum_is_ie_ie_i^\dagger$ the equation
reads $(p_a+s_i)\braket{f_a}{Xe_i}=\braket{Cf_a}{\Delta\Pi e_i}$, and
$\lVert Cf_a\rVert\leq(\lambda/2)^{-1/2}p_a$, so every entry of $X$ has modulus at most
$(\lambda/2)^{-1/2}\lVert\Delta\rVert$; the entries with $p_a=s_i=0$ vanish because
$X=PC^{-1}-C^{-1}P_\infty$. This is the Schur-product argument of the relative perturbation
bounds for the unitary polar factor \cite{Li1997Relative}. The Gram estimate of
\cite{gap:mswc24_block_form_mixed_overlap} at unit weights gives
$\lVert\Delta\rVert\leq Ke^{-2\gamma q/\xi}$, independently of the weights, so for all
$q$ beyond a threshold
\begin{align*}
    \lVert(P-P_\infty)C^{-1}\rVert\leq Ke^{-\gamma q/\xi},
\end{align*}
uniformly in the nonzero weights, which need not satisfy $\lvert\mu_{j,k}\rvert\leq1$. The
blocks $L_j^\dagger PC^{-1}L_{j'}$ are therefore $\delta_{jj'}(\sqrt{\sigma_j}^{\,T}\otimes
\mathbf 1)$ up to $Ke^{-\gamma q/\xi}$.

The row $L_j^\dagger P$, read as a tensor, is the direct sum over the copies $(j',k)$ with the
weights $\mu_{j',k}^q$ of the tensors read from the blocks $L_j^\dagger PC^{-1}L_{j'}$, so
\begin{align*}
    z_j=\sum_{j'}\beta_{j'}w_{jj'},\qquad
    w_{jj'}=\braket{\Omega_j}{\phi_M(L_j^\dagger PC^{-1}L_{j'})},
\end{align*}
and the telescoping estimate gives $\lvert w_{jj'}-\delta_{jj'}\rvert\leq C_{jj'}y\,e^{C_{jj'}y}$.
Hence
\begin{align*}
    \frac1b\sum_j\overline{\beta_j}z_j-1
    =\frac1b\sum_{j,j'}\overline{\beta_j}\beta_{j'}\bigl(w_{jj'}-\delta_{jj'}\bigr)
\end{align*}
has modulus at most $\sum_{j,j'}C_{jj'}y\,e^{C_{jj'}y}$, since
$\lvert\beta_j\beta_{j'}\rvert\leq b$, and the rest of the proof of
\cite{gap:mswc24_repeated_block_corrected_state} goes through with no division by $b^{1/2}$.

\section{What is formalized}
For blocks with multiplicities whose $q$-site states may overlap, normal in the gauge
$\sum_i(A_j^i)^\dagger A_j^i=\mathbf 1$ with positive definite fixed points $\sigma_j$ of trace
one, with $\lambda_2$ bounding the moduli of the eigenvalues other than one of every transfer
matrix of a block and of all eigenvalues of the mixed transfer matrices of distinct blocks,
and $0<\gamma<1/2$, there are $C,C'>0$ such that for all nonzero complex weights, every $q$
and every $M\geq1$ with $\beta\neq0$, the corrected approximating state satisfies
\begin{align*}
    \varepsilon\leq Cy\,e^{Cy},\qquad\varepsilon\leq C'y,\qquad y=Me^{-\gamma q/\xi}.
\end{align*}
The constants do not depend on the weights, and the normalization $\lvert\mu_{j,k}\rvert\leq1$
is not used. The same bounds hold for the source's state with the pairs on one copy of each
block and the corrected coefficients of
\cite{gap:mswc24_multiplicity_fixed_point}.\footnote{Formal declarations:
    \leanid{Matrix.norm_le_of_mul_add_mul_eq_mul} in
    \doclink{TNLean/Algebra/MatrixSylvesterBound};
    \leanid{MPSTensor.norm_polarPos_sub_mul_sum_le} and
    \leanid{MPSTensor.exists_norm_relativeRow_sub_le} in
    \doclink{TNLean/MPS/Preparation/RelativePositivePart};
    \leanid{MPSTensor.mpvOverlap_conjTranspose_copyIsometry_mul_polarPos},
    \leanid{MPSTensor.exists_norm_mpvOverlap_sub_ite_le},
    \leanid{MPSTensor.exists_approximationError_le_repeatedOverlappingBlockSum} and
    \leanid{MPSTensor.exists_approximationError_le_mul_repeatedOverlappingBlockSum} in
    \doclink{TNLean/MPS/Preparation/RepeatedOverlappingBlockWeights};
    \leanid{MPSTensor.exists_approximationError_le_repeatedOverlappingBlockSum_oneCopy} and
    \leanid{MPSTensor.exists_approximationError_le_mul_repeatedOverlappingBlockSum_oneCopy}
    in \doclink{TNLean/MPS/Preparation/OneCopyPairState}.}
The remaining deviation from the source is the rate, recorded in
\cite{gap:mswc24_block_form_mixed_overlap}.

\section{Numerical check}
A script evaluates the corrected state in arithmetic of 90 significant digits for the two random
normal blocks of bond dimension two on $d=2$ used in
\cite{gap:mswc24_repeated_block_corrected_state}, with $\lvert\lambda_2\rvert=0.850$, at
$q=5,9,13$, where $N=qM$ is odd for odd $M$.
\begin{itemize}
    \item For copy weights $(1,i,-1,-i)$ on the first block and $0.5$ on the second, $\beta_1=0$
          whenever $4\nmid N$. At $M=5$ the error is $0.427$, $7.55\cdot10^{-2}$,
          $1.51\cdot10^{-2}$, while $b=0.5^{2N}$ is $8.9\cdot10^{-16}$, $8.1\cdot10^{-28}$,
          $7.3\cdot10^{-40}$; $\varepsilon/(M\lvert\lambda_2\rvert^q)$ is $0.193$,
          $6.53\cdot10^{-2}$, $2.51\cdot10^{-2}$.
    \item For copy weights $(1,-1)$ and $0.5$, at $M=1,3,5$ the ratio
          $\varepsilon/(M\lvert\lambda_2\rvert^q)$ stays between $1.78\cdot10^{-2}$ and $0.227$
          over the three values of $q$, while $b$ ranges from $9.8\cdot10^{-4}$ down to
          $7.3\cdot10^{-40}$.
\end{itemize}
The error is controlled by $y$ alone, with no growth as $b$ decreases, as the bound states.

\bibliographystyle{alpha}
\bibliography{references}
\end{document}
